% !TEX root = ../main.tex
\chapter{\texorpdfstring{$\lambda$}{Lambda}-calcolo}
\label{chap:lambda_calculus}

Il $\lambda$-calcolo rappresenta uno dei pilastri fondamentali della teoria della computabilità. Sviluppato da 
Alonzo Church negli anni '30, esso si configura come un formalismo basato esclusivamente sulla manipolazione 
sintattica di termini per definire le funzioni calcolabili. In questo capitolo verranno introdotti i costrutti 
elementari della sua grammatica, le relazioni di riduzione che ne governano la dinamica e, infine, il problema 
della terminazione attraverso l'analisi dei termini divergenti.

\section{Tesi di Church-Turing}
Il $\lambda$-calcolo è un formalismo computazionale, ovvero un insieme di regole logico-matematiche precise utilizzate per definire 
le funzioni calcolabili. Questo sistema non è un semplice linguaggio, ma rappresenta un modello di calcolo universale: secondo la 
tesi di Church-Turing, infatti, ogni funzione intuitivamente calcolabile può essere espressa attraverso questo formalismo (o altri 
equivalenti, come le Macchine di Turing). In questo contesto, il calcolo non avviene tramite la modifica di uno stato di memoria, 
ma attraverso la riscrittura e la semplificazione di espressioni, rendendo la definizione di funzione l'elemento atomico del sistema.

\section{Grammatica del \texorpdfstring{$\lambda$}{Lambda}-calcolo}
La grammatica BNF del $\lambda$-calcolo, definita in \eqref{eq:lambda_grammar}, si basa su tre costrutti fondamentali:
le varibili, l'astarzione e l'applicazione. La grammatica definisce solo il sistema di scrittura dei termini, 
essa non assegna alcuna semantica e nessuna relazione tra gli stessi ad esclusione di quelle sintattiche.
\begin{equation}
    \begin{array}{lrll}
        \text{Terms:}   & M, N & ::= x             & \text{(Variabili)}   \\
                        &      & \mid \lambda x. M & \text{(Astrazione)}  \\
                        &      & \mid M N          & \text{(Applicazione)}
    \end{array}
    \label{eq:lambda_grammar}
\end{equation}
In ordine, i costrutti prendono il nome di variabile, astrazione e applicazione; la loro arietà è rispettivamete atomica, unaria e binaria.
L'arietà all'interno della grammatica deternima il numero di elementi richiesti per la loro formazione dei termini. È inoltre opportuno osservare 
che la definizione di queste proprietà appartiene almeno per ora al metalinguaggio, ovvero il linguaggio logico-matematico e non al $\lambda$-calcolo.
    
\section{Relazioni di calcolo e sostituzione}
Le relaioni tra termini del $\lambda$-calcolo sono cio che traforma una insieme di puramente sintattico in un effettivo formalismo computazionale. 
Alcune di esse sono necessarie 

\subsection{Variabili libere e legate}
Prima di definire formalmente le relazioni di calcolo, è necessario introdurre la natura delle variabili che compaiono all'interno dei termini. L'occorrenza 
di una variabile può essere \emph{libera} o \emph{legata},un'occorrenza della variabile $x$ si dice legata se si trova sotto il raggio d'azione di un'astrazione 
$\lambda x$, mentre si dice libera in caso contrario. L'insieme delle \emph{variabili libere} $FV(M)$ (da \emph{Free Variables}) e delle \emph{variabili legate} 
$BV(M)$ (da \emph{Bound Variables}) di un termine $M$ vengono definiti ricorsivamente sulla struttura del termine come segue.
\begin{equation}
    \begin{array}{lcl}
        FV(x)            & = & \{x\}                 \\
        FV(\lambda x. M) & = & FV(M) \setminus \{x\} \\
        FV(M N)          & = & FV(M) \cup FV(N)
    \end{array}
    \qquad
    \begin{array}{lcl}
        BV(x)            & = & \emptyset        \\
        BV(\lambda x. M) & = & BV(M) \cup \{x\} \\
        BV(M N)          & = & BV(M) \cup BV(N)
    \end{array}
    \label{eq:fv_bv_definitions}
\end{equation}
Un termine $M$ si dice \emph{chiuso} (o combinatore) se non contiene variabili libere, ovvero se $FV(M) = \emptyset$.

\subsection{Sostituzione}
L'operazione di sostituzione formale è fondamentale per la $\beta$-riduzione. La notazione $M[N/x]$ (o $M[x := N]$) 
indica il termine ottenuto sostituendo l'espressione $N$ a ogni occorrenza libera della variabile $x$ nel termine $M$. 
L'operazione viene definita ricorsivamente sulla struttura del termine $M$ come segue.
\begin{equation}
    M[N/x] = 
    \begin{cases} 
        N & \text{if } M = x                        \\
        y & \text{if } M = y \text{ e } y \neq x    \\
        (P[N/x] \, Q[N/x]) & \text{if } M = (P Q)   \\
        \lambda x. P & \text{if } M = \lambda x. P  \\
        \lambda y. (P[N/x]) & \text{if } 
            \begin{cases} 
            M = \lambda y. P    \\
            y \neq x            \\
            y \notin FV(N)
            \end{cases}         \\
        \perp & \text{otherwise}
    \end{cases}
    \label{eq:substitution_algorithm}
\end{equation}
L'ultima clausola valida introduce una condizione cruciale ($y \notin FV(N)$) per evitare il fenomeno del \emph{variable capture} (cattura della variabile).
Se $y$ facesse parte delle variabili libere di $N$, il legame introdotto da $\lambda y$ ne altererebbe il significato originale dopo la sostituzione. In questo 
caso, si rende necessario applicare preventivamente la $\alpha$-equivalenza per rinominare la variabile legata $y$ con un nome nuovo.

\subsection{\texorpdfstring{$\alpha$}{alpha}-equivalenza}
la $\alpha$-equivalenza stabilisce che il nome delle variabili legate è irrilevante ai fini del significato del ternine. Due termini sono considerati 
identici se l'uno puo essere ottenuto dal'altro tramite la ridenominazione coerente delle proprie variabili legate. Questa relazione è fondametare per 
evitare il fenomeno del varible capture durante le operazione di sostituzione, garantendo che i nomi scelti arbitrariamente non interferiscano con 
le varibili libere.
\begin{equation}
    \lambda x. M \equiv_\alpha \lambda y. M[y/x] \quad \text{se } y \notin FV(M)
    \label{eq:alpha_conversion}
\end{equation}

\subsection{\texorpdfstring{$\beta$}{beta}-riduzione}
\label{sec:beta_reduction}
La $\beta$-riduzione stabilisce il motore del $\lambda$-calcolo e formalizza il cncetto di applicazione funzionale. Consiste nel sostituire l'argomento N a ogni
occorrenza libera della variabile x all'interno del corpo della funzione M.
\begin{equation}
    (\lambda x. M) N \to_\beta M[x := N]
    \label{eq:beta_reduction}
\end{equation}
Un ternine che presenta questa struttura si dice \emph{redex} (reducible expression). Il calcolo consiste dell'applicazione ripetura di questa regola, 
quando un termine non contiene piu redex si dice che è in \emph{forma normale}. La forma normale è il risultato concreto della computazione.

\subsection{\texorpdfstring{$\eta$}{eta}-conversione}
La $\eta$-conversione ($\to_\eta$) introduce il principio di estensionalità, secondo il quale due funzioni sono uguali se, 
applicate allo stesso argomento, producono il medesimo risultato.
\begin{equation}
    \lambda x. (M x) \to_\eta M \quad \text{se } x \notin FV(M)
    \label{eq:eta_conversion}
\end{equation}
Questa operazione permette di semplificare espressioni ridondanti, eliminando astrazioni che si limitano a "passare" 
l'argomento a un altro termine.

\section{Problemi di divergenza}
Nel $\lambda$-calcolo, come esposto nella sezione \ref{sec:beta_reduction}, la computazione deriva dalla riduzione dei termini. Quando questa non è possibile allora
il termine è in forma normale, ovvero il risultato della computazione. Tuttavia è importante osservare che non tutti i termini ammetono una forma normale, in alcuni 
casi la riduzione può proseguire indefinitivamente. 

\subsection{Definizione di divergenza}
L'assenza di una forma normale è ciò che definisce formalmente la \emph{divergenza}. 
In un sistema Turing-completo come il $\lambda$-calcolo non tipato, la divergenza è 
l'equivalente funzionale di un ciclo infinito (loop) in un linguaggio di programmazione imperativo.

\subsection{Il termine \texorpdfstring{$\Omega$}{Omega}}
L'esempio piu semplice di termine divergente è il combinatore $\Omega$. Esso viene costruito a partire dal termine di auto-applicazione $\omega$ 
(talvolta indicato come $\delta$).
\begin{equation}
    \omega = \lambda x. (x x)
\end{equation}
Dalla regola di $\beta$-riduzione appare evidente che l'auto-applicazione del termine $\omega$ sia divergente, la contrazione del redex non conduce a una forma 
normale ma riporta al termine iniziale.